% chapters/17_prompting.tex
% Part of: AI / LLM / Agents / Hardware Notes

\section{Prompt Engineering \& In-context Learning}

\subsection{Zero-shot \& few-shot prompting}

\textbf{Zero-shot prompting} asks the model to perform a task without any examples,
relying entirely on pretrained knowledge. Effective for clear, well-defined tasks where
the instruction alone is sufficient.

\textbf{Few-shot prompting} provides $k$ (typically 1--8) exemplar input/output pairs
in the prompt before the actual query. The model infers the pattern and generalises to
new inputs without weight updates. Key properties:

\begin{itemize}
  \item Demonstrated format strongly constrains output structure — useful for JSON,
    tables, and constrained classification.
  \item Example \emph{quality} matters more than quantity; misleading exemplars degrade
    performance.
  \item Sensitive to example ordering (recency bias); shuffling exemplars can improve
    robustness.
  \item In LangChain, \texttt{FewShotPromptTemplate} formalises this pattern,
    composing a \texttt{PromptTemplate} with stored exemplar pairs and an
    \texttt{ExampleSelector} for dynamic retrieval from a large pool.
\end{itemize}

\subsection{Chain-of-thought (CoT)}

Chain-of-thought prompting encourages the model to emit a step-by-step reasoning trace
before its final answer, substantially improving performance on arithmetic, commonsense,
and multi-step symbolic reasoning tasks.

\begin{itemize}
  \item \textbf{Zero-shot CoT:} append ``\emph{Let's think step by step}'' (or
    equivalent) to the prompt. Requires no worked examples.
  \item \textbf{Few-shot CoT:} provide exemplar triples of (question, reasoning chain,
    answer) so the model imitates the reasoning style and level of detail.
  \item \textbf{Emergent property:} CoT is effective only above roughly 100B parameters
    in base models; smaller models produce incoherent rationales that hurt accuracy.
    Instruction-tuning (e.g.\ IBM Granite Instruct, GPT-4o-mini) transfers CoT
    capability to much smaller models via distillation on CoT traces.
\end{itemize}

\textbf{CoT vs.\ prompt chaining --- a critical distinction:}
\begin{itemize}
  \item \emph{CoT} elicits multi-step reasoning \emph{within a single prompt turn};
    reasoning lives in the model's output for that turn.
  \item \emph{Prompt chaining} sequences multiple API calls, where the output of one
    call becomes context for the next --- a structural/orchestration pattern, not a
    reasoning one.
\end{itemize}

\textbf{Extended thinking} (Claude) extends CoT with a dedicated scratchpad block
that is invisible in the conversation history, enabling deeper reasoning without
consuming visible response tokens. Each thinking block carries a cryptographic
signature for tamper detection in multi-turn loops (see Chapter~52).

\begin{paperbox}[title=ReAct --- CoT meets tool use (Yao et al.\, 2023)]
  ReAct integrates chain-of-thought \emph{Thought} steps with \emph{Action} steps
  (tool invocations) and \emph{Observation} steps (tool results) in an alternating
  loop. This grounds reasoning in real-world evidence and reduces hallucination.
  Full coverage in Chapter~23.
\end{paperbox}

\subsection{Structured output prompting}

Agentic pipelines require machine-parseable outputs. A tool call with a malformed
JSON argument crashes the loop; a stray natural-language prefix breaks a downstream
parser. Structured output is therefore load-bearing, not cosmetic.

\textbf{Techniques (roughly increasing reliability):}
\begin{itemize}
  \item Describe the target schema in natural language within the system prompt.
  \item Provide a few-shot example of the exact JSON/XML format expected.
  \item Use API-level JSON mode (OpenAI) or tool schemas (Anthropic) to constrain the
    model's decoding to syntactically valid output.
  \item Use \texttt{PromptTemplate} classes (LangChain, LlamaIndex) that bake the
    format requirement into a reusable, parameterised template object.
\end{itemize}

\textbf{Tool schema design} (see also Chapter~20) defines the structured interface
the model uses to call external functions: unique name, natural-language description,
parameter names/types, required vs.\ optional fields. The description is the
primary signal the model uses to decide \emph{when} to invoke the tool.

\subsection{System prompts \& jailbreaks}

The \textbf{system prompt} establishes the model's role, constraints, and behavioural
guidelines before the conversation begins. It is the primary mechanism for customising
base model behaviour without fine-tuning.

\textbf{In agentic systems the system prompt typically:}
\begin{itemize}
  \item Defines the agent's persona, available tools, and expected output format.
  \item Lists prohibited actions and safety rules (``\emph{in-context}'' controls).
  \item Embeds Skills documentation (e.g.\ Cowork's SKILL.md files) encoding task-specific
    best practices that fire reliably for specific requests.
  \item Is cached at the API level (Anthropic prompt caching, OpenAI prefix caching)
    so it does not re-incur compute cost on every turn.
\end{itemize}

\textbf{The fundamental weakness of system-prompt safety:} in-context instructions
only constrain the model if it ``chooses'' to follow them. Sufficiently adversarial
inputs --- prompt injection via observed web/email content, context overflow, or
multi-step reasoning-around-the-rules --- can cause the model to ignore them.
This is why out-of-process runtime enforcement (OpenShell, VM sandboxing) is needed
as a backstop (Chapter~42).

\begin{warnbox}[title=Jailbreaks]
  Jailbreak techniques include: role-play framing (``pretend you have no restrictions''),
  token manipulation (unusual Unicode, homoglyphs), base64 encoding of forbidden
  requests, prompt injection via tool results, and suffix-appending adversarial
  attacks. Defences: constitutional AI (Chapter~43), input/output classifiers,
  and runtime enforcement (Chapter~42).
\end{warnbox}

